\section{Penelitian Terkait}

\subsection{Penelitian IndoBERT pada Analisis Sentimen}

IndoBERT merupakan \textit{Pre-trained Language Model} (PLM) berbasis arsitektur Transformer yang dirancang khusus untuk bahasa Indonesia. Model ini banyak digunakan pada tugas analisis sentimen karena mampu memahami hubungan semantik dan konteks kalimat secara lebih baik dibandingkan metode tradisional. Penelitian Yulianti dan Nissa \cite{yulianti2023} menunjukkan bahwa IndoBERT memberikan performa yang baik pada tugas \textit{Aspect-Based Sentiment Analysis} (ABSA), baik menggunakan pendekatan \textit{single-sentence} maupun \textit{sentence-pair classification}. Hasil serupa juga ditemukan oleh Latifah dkk. \cite{latifah2024} dan Wafda dkk. \cite{wafda2025} yang menunjukkan efektivitas IndoBERT dalam mengklasifikasikan sentimen pada berbagai domain berbahasa Indonesia.

Selain itu, Apriliani \cite{apriliani2025} menunjukkan bahwa \textit{fine-tuning} IndoBERT pada ulasan hotel berbahasa Indonesia mampu menghasilkan performa klasifikasi sentimen yang sangat baik. Mustofa dkk. \cite{mustofa2026} juga memanfaatkan model BERT untuk analisis sentimen berbasis aspek pada aplikasi E-SKM Jawa Tengah dan menunjukkan bahwa model berbasis Transformer mampu meningkatkan kualitas identifikasi sentimen dan aspek dibandingkan pendekatan konvensional.

\subsection{Penelitian BERTopic pada Pemodelan Topik}

BERTopic merupakan metode pemodelan topik modern yang memanfaatkan representasi \textit{embedding} berbasis Transformer untuk menghasilkan topik yang lebih koheren dibandingkan metode tradisional seperti LDA dan NMF. Metode ini mengombinasikan \textit{sentence embedding}, reduksi dimensi menggunakan UMAP, dan klasterisasi menggunakan HDBSCAN untuk menemukan topik secara otomatis dari kumpulan dokumen.

Penelitian Babalola dkk. \cite{babalola2023} menunjukkan bahwa BERTopic menghasilkan topik yang lebih representatif dan mudah diinterpretasikan dibandingkan LDA dan NMF. Temuan tersebut diperkuat oleh Janssens dkk. \cite{janssens2025} yang menyatakan bahwa BERTopic mampu mempertahankan kualitas topik meskipun dilakukan proses reduksi topik sehingga meningkatkan interpretabilitas hasil pemodelan.

\subsection{Penelitian Aspect-Based Sentiment Analysis pada Pariwisata}

Penerapan \textit{Aspect-Based Sentiment Analysis} (ABSA) pada bidang pariwisata berkembang seiring meningkatnya jumlah ulasan wisatawan pada platform digital. Afendi dkk. \cite{afendi2023} menjelaskan bahwa pemanfaatan \textit{Pre-trained Language Model} pada ekosistem \textit{smart tourism} mampu membantu pengelola memahami persepsi wisatawan melalui analisis ulasan secara otomatis.

Penelitian Ramadhani dkk. \cite{ramadhani2024} berhasil memanfaatkan IndoBERT untuk menganalisis persepsi wisatawan berdasarkan data Google Reviews pada destinasi wisata di Bandung Raya. Selain itu, Nalury dan Irawan \cite{nalury2024} menerapkan analisis sentimen berbasis IndoBERT pada wisata religi di Pulau Jawa, sedangkan Al-Hafiz dkk. \cite{alhafiz2024} mengintegrasikan IndoBERT dan BERTopic untuk menganalisis ulasan objek wisata di Yogyakarta sehingga menghasilkan identifikasi aspek dan sentimen yang lebih rinci.

\subsection{Penelitian Integrasi IndoBERT dan BERTopic}

Integrasi IndoBERT dan BERTopic mulai banyak digunakan untuk menggabungkan analisis sentimen dan pemodelan topik secara bersamaan. Kurniawan dkk. \cite{kurniawan2023} menunjukkan bahwa kombinasi kedua metode tersebut mampu menghasilkan pemetaan aspek dan sentimen yang lebih rinci pada ulasan pariwisata Indonesia. Hasil serupa juga diperoleh oleh Girsang dkk. \cite{girsang2024} dan Maylawati dkk. \cite{maylawati2024} yang menunjukkan bahwa integrasi IndoBERT dan BERTopic mampu meningkatkan kualitas identifikasi aspek dan sentimen.

Selain pada domain pariwisata, pendekatan serupa juga diterapkan oleh Nur dkk. \cite{nur2025} pada analisis merek teknologi di media sosial dan oleh Sihombing dan Widiyaningtyas \cite{sihombing2025} pada isu kesehatan masyarakat. Hasil penelitian tersebut menunjukkan bahwa kombinasi model berbasis Transformer dan pemodelan topik modern mampu menghasilkan informasi yang lebih komprehensif dibandingkan penggunaan salah satu metode secara terpisah.
